\section{Systemüberblick}
\label{sec:systemueberblick}

StudyMummy ist eine webbasierte Lernplattform, die hochgeladene Lernunterlagen
in Aufgaben, Quizfragen und Merkblätter überführt und Lernende in einem
sokratischen Chat adaptiv unterstützt. Der agentische Kern verarbeitet pro Chat-Turn eine
authentifizierte Nachricht zusammen mit Aufgabe, Hilfestufe, Dialoghistorie und
abgerufenem Dokumentwissen. Planner, Tutor und Reviewer besitzen getrennte
Ziele, Befugnisse und laufbezogene Zustände. Über typisierte Nachrichten kann
der Reviewer einen Entwurf freigeben, eine Tutorrevision verlangen oder eine
begrenzte Neuplanung anstoßen. Der Orchestrator verteilt diese Nachrichten und
begrenzt den Ablauf, trifft aber keine fachliche Entscheidung.

Technisch besteht das System aus einer Angular-22-Single-Page-Application,
einer asynchronen FastAPI-Anwendung mit Pydantic und SQLAlchemy sowie
PostgreSQL~15 mit \texttt{pgvector}. Ein OpenAI-kompatibler Adapter stellt
strukturierte Aufrufe großer Sprachmodelle (Large Language Models, LLMs) und
Function Calling bereit. Docker Compose trennt Entwicklungs- und
Produktionsprofile. Im Produktivprofil liefert Nginx das
Frontend aus und leitet API- und WebSocket-Verkehr an das Backend weiter.
Persistiert werden unter anderem Nutzer, Dokumente, Vektorchunks, Aufgaben,
Sitzungen, Chatlogs, Lernprofile und Gamification-Zustände. Die zentrale Frage
lautet: \emph{\Forschungsfrage}
